\subsection{算子代数}
\label{sec:ch4-operator-algebra}

设 $P(\xi)=\sum_\alpha b_\alpha\xi^\alpha$ 是 $n$ 个变量的常系数多项式, $P(D)$ 是相应的微分多项式, 即
\[
 P(D)f=\sum_\alpha b_\alpha D^\alpha f.
\]
由 Fourier 变换的微分性质,
\[
 \widehat{P(D)f}(\xi)=P(2\pi i\xi)\widehat f(\xi).
\]
定义 $\Lambda$ 为正算子 $-\Delta$ 的平方根:
\begin{equation}
\label{eq:ch4-lambda-definition}
 \widehat{\Lambda f}(\xi)=2\pi|\xi|\widehat f(\xi).
\end{equation}
若 $P$ 是 $m$ 次齐次多项式, 则
\begin{equation}
\label{eq:ch4-differential-operator-factorization}
 P(D)f=T(\Lambda^mf),
\end{equation}
其中 $T$ 定义为
\[
 \widehat{Tf}(\xi)=i^m\frac{P(\xi)}{|\xi|^m}\widehat f(\xi).
\]

\hypertarget{chapter-04-page-093}{}
乘子 $P(\xi)/|\xi|^m$ 是零次齐次的, 在 $S^{n-1}$ 上的限制等于 $P(\xi')$, 因而是 $C^\infty$ 函数. 事实上, 可以证明 $T$ 是若干 Riesz 变换复合的线性组合. $T$ 不一定是 \eqref{eq:ch4-singular-integral-definition} 型的奇异积分, 因为 $P(\xi)$ 在 $S^{n-1}$ 上的平均值未必为零. 但减去一个常数, 即恒等算子的常数倍, 就可得到这样的奇异积分.

\begin{Theorem}
\label{thm:ch4-smooth-homogeneous-multiplier-representation}
若 $m\in C^\infty(\mathbb R^n\setminus\{0\})$ 是零次齐次函数, 且 $T_m$ 由 $\widehat{T_mf}=m\widehat f$ 定义, 则存在 $a\in\mathbb C$ 以及平均值为零的 $\Omega\in C^\infty(S^{n-1})$, 使得对每个 $f\in\mathcal S$,
\[
 T_mf=af+\operatorname{p.v.}\frac{\Omega(x')}{|x|^n}*f.
\]
\end{Theorem}

任何零次齐次函数都是一个常数与一个在 $S^{n-1}$ 上平均值为零的零次齐次函数之和, 因而定理~\ref{thm:ch4-smooth-homogeneous-multiplier-representation} 是下面引理的直接推论.

\begin{Lemma}
\label{lem:ch4-zero-mean-multiplier-kernel}
设 $m\in C^\infty(\mathbb R^n\setminus\{0\})$ 是在 $S^{n-1}$ 上平均值为零的零次齐次函数. 则存在平均值为零的 $\Omega\in C^\infty(S^{n-1})$, 使得
\[
 \widehat m(\xi)=\left(\operatorname{p.v.}
 \frac{\Omega(x')}{|x|^n}\right)(\xi).
\]
\end{Lemma}

\begin{Proof}
把 $m$ 看作缓增分布. 对每个 $i$, $\partial^nm/\partial x_i^n$ 在 $\mathbb R^n\setminus\{0\}$ 上光滑, 是 $-n$ 次齐次的, 且在 $S^{n-1}$ 上平均值为零. 又因为普通分布导数与其主值延拓之差支于原点,
\[
 \frac{\partial^nm}{\partial x_i^n}
 =\operatorname{p.v.}\frac{\partial^nm}{\partial x_i^n}
 +\sum_{|\alpha|\le k}C_\alpha D^\alpha\delta.
\]
取 Fourier 变换后, 左端及右端第一项都是零次齐次分布, 因而右端的多项式必须退化为常数. 所以在 $\mathbb R^n\setminus\{0\}$ 上, $\widehat m$ 与一个 $-n$ 次齐次光滑函数一致. 记其在 $S^{n-1}$ 上的限制为 $\Omega$.

取支于环域 $1<|x|<2$ 的正径向 Schwartz 函数 $\varphi$. 由径向性, 齐次性以及 $m$ 在球面上的平均值为零可得 $\widehat m(\varphi)=0$, 从而 $\Omega$ 的球面平均值也为零. 最后, $\widehat m-\operatorname{p.v.}\Omega(x')/|x|^n$ 支于原点. 其 Fourier 反变换是多项式; 因两项的 Fourier 反变换都有界, 该多项式只能是常数. 两者的球面平均值都为零, 所以这个常数为零.
\end{Proof}

\hypertarget{chapter-04-page-094}{}
\begin{Theorem}
\label{thm:ch4-singular-integral-operator-algebra}
由定理~\ref{thm:ch4-smooth-homogeneous-multiplier-representation} 定义的算子全体 $\mathcal A$ 构成交换代数. $\mathcal A$ 中的元素可逆, 当且仅当其乘子 $m$ 在 $S^{n-1}$ 上处处不为零.
\end{Theorem}

\begin{Proof}
若 $m_1,m_2$ 对应于 $T_{m_1},T_{m_2}$, 则 $T_{m_1}T_{m_2}=T_{m_1m_2}$, 所以 $\mathcal A$ 是交换代数. 恒等算子的乘子是 $1$, 因而 $T_m$ 可逆当且仅当 $T_{1/m}\in\mathcal A$. 而 $1/m\in C^\infty(\mathbb R^n\setminus\{0\})$ 当且仅当 $m(\xi)\ne0$ 对每个 $\xi\in S^{n-1}$ 成立.
\end{Proof}

\eqref{eq:ch4-differential-operator-factorization} 中与多项式 $P$ 相应的算子 $T$ 可逆, 当且仅当 $P(\xi)$ 在 $S^{n-1}$ 上处处不为零, 即 $P(D)$ 是椭圆算子. 若 $P$ 的系数为实数, 则 $m$ 必为偶数, 且 $\Lambda^m=(-\Delta)^{m/2}$. 因而求解 $P(D)u=f$ 的问题化为求解 $(-\Delta)^{m/2}u=T^{-1}f$. 关于这个算子的更多内容见第~\ref{subsec:ch4-fractional-integrals-sobolev}~节.
